\scanpage{13}
\textcircled{2} Résultats auxiliaires sur les relations d'équivalence

\uline{Proposition 1} \quad Soient $A$ un anneau, $\unsure{m|\mathfrak m}$ un idéal maximal de $A$, $A'$ une $A$-algèbre, contenant $A$, tel que $A'/A$ soit un $A$-module isomorphe à $A/m=k$

Sous ces conditions

(i) Il \deleted{\unsure{n'}}y a dans $\operatorname{Spec}A'$ sur $\operatorname{Spec}A$ exactement deux relations d'\unsure{équivalences|équivalence}, la plus fine \deleted{et} [correspondant à la diagonale de $\operatorname{Spec}A'$ sur $\operatorname{Spec} A$] et la moins fine [correspondant au produit fibré]

(ii) la suite
\[
\deleted{\unsure{0\longrightarrow}}A\longrightarrow A'\rightrightarrows A'\otimes_A A'
\]
est exacte, \deleted{\unsure{plus pr.|plus précis}} i.e. $H^0(A'/A)\xleftarrow{\sim}A$, et $H^i(A'/A)=0$ pour $i\ne0$

(iii) $\operatorname{Spec}A'\longrightarrow\operatorname{Spec}A$ est un morphisme de descente effective pour la catégorie fibrée des Modules \unsure{qu. coh.:|qu. coh.} plats, et reste \unsure{telle|telles|tel} pour \unsure{toute|toutes} \unsure{extension|extensions} de la base $X\longrightarrow\operatorname{Spec}A$).
% archive sheet number: 12
